\section{The Zurich Tradition: Twenty Papers by Ernst Fehr (2010--2025)}
\label{app:fehrpapers}
%%%%%%%%%%%%%%%%%%%%%%%%%%%%%%%%%%%%%%%%%%%%%%%%%%%%%%%%%%%%%%%%%%%%%%%%%%%%%%%

This appendix integrates twenty papers by Ernst Fehr and collaborators 
published since 2010, focusing on behavioral economics contributions
(excluding pure neuroeconomics). This demonstrates the framework's
ability to systematically organize the research program of a leading
behavioral economist---and, not coincidentally, connects to the
intellectual tradition from which this framework emerged.

\subsection{Overview: Fehr's Research Program in Framework Terms}

Ernst Fehr's research since 2010 addresses several core framework components:

\begin{center}
\renewcommand{\arraystretch}{1.2}
\small
\begin{tabular}{ll}
\hline
\textbf{Framework Element} & \textbf{Fehr's Contribution} \\
\hline
$C^{SS}_{ij} \neq 0$ & Social preferences: inequality aversion, reciprocity \\
$U \neq U^F$ only & Other-regarding preferences enter utility \\
$\Psi_{norm}$ & Social norms shape behavior \\
$\Psi_{inst}^{culture}$ & Corporate/banking culture affects decisions \\
Heterogeneity in $C$ & Distribution of preference types \\
$C^*$ stability & Cooperation and punishment dynamics \\
\hline
\end{tabular}
\end{center}

\subsection{Twenty Papers Integrated}

%--- Paper 1 ---
\subsubsection{Fehr \& Charness (2025): Social Preferences Survey}

\paragraph{Citation.}
Fehr, E. \& Charness, G. (2025). ``Social Preferences: Fundamental 
Characteristics and Economic Consequences.'' \emph{Journal of Economic Literature}.

\paragraph{Core Contribution.}
Comprehensive survey establishing that three preference types dominate:
altruistic, inequality averse, and predominantly selfish---with inequality
aversion being the majority.

\paragraph{Framework Integration.}
\begin{itemize}
  \item \textbf{Heterogeneity:} Population distribution of $(\alpha, \beta)$ 
        in Fehr-Schmidt utility
  \item \textbf{Clustering:} Bayesian methods reveal discrete types, not continuum
  \item \textbf{Prediction:} Type distribution predicts cooperation, redistribution
\end{itemize}

%--- Paper 2 ---
\subsubsection{Berger, Fehr et al. (2025): Working Memory Training}

\paragraph{Citation.}
Berger, E.M., Fehr, E., Hermes, H., Schunk, D., \& Winkel, K. (2025).
``The Impact of Working Memory Training on Children's Cognitive and Noncognitive Skills.''
\emph{Journal of Political Economy}.

\paragraph{Framework Integration.}
\begin{itemize}
  \item \textbf{Self-regulation:} Affects future $U^F$, $U^P$ outcomes
  \item \textbf{Early intervention:} Modifying $\Psi_{individual}$ early has lasting effects
  \item \textbf{Cross-domain:} Cognitive training $\to$ noncognitive skills $\to$ life outcomes
\end{itemize}

%--- Paper 3 ---
\subsubsection{Epper, Fehr et al. (2024): Inequality Aversion and Redistribution}

\paragraph{Citation.}
Epper, T., Fehr, E., Kreiner, C.T., Leth-Petersen, S., et al. (2024).
``Inequality Aversion Predicts Support for Public and Private Redistribution.''
\emph{PNAS}.

\paragraph{Framework Integration.}
\begin{itemize}
  \item \textbf{Preferences $\to$ Politics:} $\alpha, \beta$ parameters predict voting
  \item \textbf{Cross-domain:} $C^{SS} \to C^{Political}$---social preferences shape policy
  \item \textbf{Redistribution:} $U^{KNU}$ (collective utility) enters decisions
\end{itemize}

%--- Paper 4 ---
\subsubsection{Efferson, Fehr et al. (2024): Super-Additive Cooperation}

\paragraph{Citation.}
Efferson, C., Bernhard, H., Fischbacher, U., \& Fehr, E. (2024).
``Super-additive Cooperation.'' \emph{Nature}.

\paragraph{Core Finding.}
Combination of repeated interactions AND group competition produces
cooperation rates exceeding what either mechanism achieves alone.

\paragraph{Framework Integration.}
\begin{itemize}
  \item \textbf{Complementarity in mechanisms:} $C^{repetition, competition} > 0$
  \item \textbf{Super-additivity:} Effect of combined mechanisms exceeds sum
  \item \textbf{Context interaction:} $\Psi_{repeated} \times \Psi_{competition}$
\end{itemize}

\paragraph{Framework Insight.}
This is a direct demonstration of complementarity: mechanisms are
complements, not substitutes. The framework predicts such interactions.

%--- Paper 5 ---
\subsubsection{Epper, Fehr et al. (2024): Social Preferences Across Subject Pools}

\paragraph{Citation.}
Epper, T., Senn, J., \& Fehr, E. (2024). ``Social Preferences Across Subject Pools.''
\emph{Working Paper}.

\paragraph{Framework Integration.}
\begin{itemize}
  \item \textbf{External validity:} Do lab results generalize?
  \item \textbf{$\Psi_{sample}$:} Student vs. population samples differ
  \item \textbf{Robustness:} Core type structure holds across pools
\end{itemize}

%--- Paper 6 ---
\subsubsection{Henkel, Fehr et al. (2024): Beliefs About Inequality}

\paragraph{Citation.}
Henkel, A., Fehr, E., Senn, J., \& Epper, T. (2025).
``Beliefs About Inequality and the Nature of Support for Redistribution.''
\emph{Journal of Public Economics}.

\paragraph{Framework Integration.}
\begin{itemize}
  \item \textbf{Beliefs matter:} $\Psi_{beliefs}$ about inequality shapes policy support
  \item \textbf{Not just preferences:} Both $(\alpha, \beta)$ AND beliefs predict behavior
  \item \textbf{Information effects:} Correcting beliefs changes support
\end{itemize}

%--- Paper 7 ---
\subsubsection{Epper, Fehr et al. (2022): Preferences Predict Crime}

\paragraph{Citation.}
Epper, T., Fehr, E., Hvidberg, K.B., Kreiner, C.T., et al. (2022).
``Preferences Predict Who Commits Crime Among Young Men.'' \emph{PNAS}.

\paragraph{Framework Integration.}
\begin{itemize}
  \item \textbf{Time preferences:} Low patience predicts criminal behavior
  \item \textbf{Cross-domain prediction:} Lab measure $\to$ real-world outcome
  \item \textbf{$U^F$ vs. $U^{long-term}$:} Discount rate $\delta$ matters
\end{itemize}

%--- Paper 8 ---
\subsubsection{Schunk, Fehr et al. (2022): Teaching Self-Regulation}

\paragraph{Citation.}
Schunk, D., Berger, E., Hermes, H., Winkel, K., \& Fehr, E. (2022).
``Teaching Self-Regulation.'' \emph{Nature Human Behaviour}.

\paragraph{Framework Integration.}
\begin{itemize}
  \item \textbf{Intervention:} $\Delta\Psi_{self-regulation} \to \Delta$ outcomes
  \item \textbf{Scalable:} Low-cost classroom intervention
  \item \textbf{Persistence:} Effects last beyond intervention period
\end{itemize}

%--- Paper 9 ---
\subsubsection{Bartling, Fehr et al. (2021/2025): Trust and Contract Enforcement}

\paragraph{Citation.}
Bartling, B., Fehr, E., Huffman, D., \& Netzer, N. (forthcoming).
``The Complementarity Between Trust and Contract Enforcement.''
\emph{Economic Journal}.

\paragraph{Core Finding.}
Trust and formal contracts are \emph{complements}, not substitutes.
Better contract enforcement increases trust.

\paragraph{Framework Integration.}
\begin{itemize}
  \item \textbf{Direct complementarity:} $C^{trust, contract} > 0$
  \item \textbf{Contra crowding-out:} Institutions don't destroy intrinsic motivation
  \item \textbf{$\Psi_{inst}$ and $\Psi_{norm}$:} Reinforce each other
\end{itemize}

\paragraph{Framework Significance.}
This directly tests a core framework claim: institutional and normative
context can be complements.

%--- Paper 10 ---
\subsubsection{Bartling, Fehr et al. (2021): Markets and Moral Values}

\paragraph{Citation.}
Bartling, B., Fehr, E., \& \"Ozdemir, Y. (2021).
``Does Market Interaction Erode Moral Values?''
\emph{Review of Economics and Statistics}.

\paragraph{Framework Integration.}
\begin{itemize}
  \item \textbf{$\Psi_{market}$ effect:} Does market exposure change preferences?
  \item \textbf{Endogenous norms:} $\Psi_{norm,t+1} = f(\Psi_{norm,t}, market_t)$?
  \item \textbf{Finding:} Limited evidence for erosion---norms are sticky
\end{itemize}

%--- Paper 11 ---
\subsubsection{Fehr, Powell \& Wilkening (2021): Behavioral Implementation}

\paragraph{Citation.}
Fehr, E., Powell, M., \& Wilkening, T. (2021).
``Behavioral Constraints on the Design of Subgame-Perfect Implementation Mechanisms.''
\emph{American Economic Review}.

\paragraph{Framework Integration.}
\begin{itemize}
  \item \textbf{Mechanism design:} Standard theory assumes $C_{ij} = 0$
  \item \textbf{Reality:} Social preferences ($C_{ij} \neq 0$) constrain mechanisms
  \item \textbf{Implication:} Optimal mechanisms must account for behavioral types
\end{itemize}

%--- Paper 12 ---
\subsubsection{Al\'os-Ferrer, Fehr \& Netzer (2021): Recovering Preferences}

\paragraph{Citation.}
Al\'os-Ferrer, C., Fehr, E., \& Netzer, N. (2021).
``Time Will Tell: Recovering Preferences When Choices Are Noisy.''
\emph{Journal of Political Economy}.

\paragraph{Framework Integration.}
\begin{itemize}
  \item \textbf{Noise in decisions:} $C^{observed} \neq C^{true}$
  \item \textbf{Methodology:} Repeated observations reveal true preferences
  \item \textbf{Implication:} Single observations may misclassify types
\end{itemize}

%--- Paper 13 ---
\subsubsection{Epper, Fehr et al. (2020): Time Discounting and Wealth}

\paragraph{Citation.}
Epper, T., Fehr, E., Fehr-Duda, H., et al. (2020).
``Time Discounting and Wealth Inequality.''
\emph{American Economic Review}.

\paragraph{Core Finding.}
Patience measured in lab predicts wealth accumulation over 15 years.
Explains significant portion of wealth inequality.

\paragraph{Framework Integration.}
\begin{itemize}
  \item \textbf{Discount rate:} $\delta$ in utility affects $U^F_{long-term}$
  \item \textbf{Predictive power:} Lab measure $\to$ real wealth
  \item \textbf{Inequality:} Preference heterogeneity $\to$ outcome inequality
\end{itemize}

%--- Paper 14 ---
\subsubsection{Declerck, Fehr et al. (2020): Oxytocin and Trust (Replication)}

\paragraph{Citation.}
Declerck, C.H., Boone, C., Pauwels, L., Vogt, B., \& Fehr, E. (2020).
``A Registered Replication Study on Oxytocin and Trust.''
\emph{Nature Human Behaviour}.

\paragraph{Framework Integration.}
\begin{itemize}
  \item \textbf{Biological context:} Hormones as part of $\Psi_{biological}$
  \item \textbf{Replication:} Tests robustness of prior findings
  \item \textbf{Methodological:} Importance of pre-registration
\end{itemize}

%--- Paper 15 ---
\subsubsection{Bruhin, Fehr \& Schunk (2019): Distribution of Social Preferences}

\paragraph{Citation.}
Bruhin, A., Fehr, E., \& Schunk, D. (2019).
``The Many Faces of Human Sociality.''
\emph{Journal of the European Economic Association}.

\paragraph{Framework Integration.}
\begin{itemize}
  \item \textbf{Type distribution:} Empirical clustering of preference types
  \item \textbf{Stability:} Types stable across time and contexts
  \item \textbf{Heterogeneity:} Population is mixture, not representative agent
\end{itemize}

%--- Paper 16 ---
\subsubsection{Efferson, Fehr \& Vogt (2019): Social Influence and Traditions}

\paragraph{Citation.}
Efferson, C., Fehr, E., \& Vogt, S. (2019).
``The Promise and Peril of Using Social Influence to Reverse Harmful Traditions.''
\emph{Nature Human Behaviour}.

\paragraph{Framework Integration.}
\begin{itemize}
  \item \textbf{Norm change:} Policy to shift $\Psi_{norm}$
  \item \textbf{Tipping points:} Critical mass for norm transition
  \item \textbf{Risk:} Interventions can backfire if poorly designed
\end{itemize}

%--- Paper 17 ---
\subsubsection{Engelmann, Fehr et al. (2019): Antisocial Personality}

\paragraph{Citation.}
Engelmann, J.B., Fehr, E., Schmid, B., et al. (2019).
``On the Psychology and Economics of Antisocial Personality.''
\emph{PNAS}.

\paragraph{Framework Integration.}
\begin{itemize}
  \item \textbf{Extreme types:} Antisocial as tail of distribution
  \item \textbf{Cross-validation:} Psychological and economic measures align
  \item \textbf{Policy:} Identifies intervention targets
\end{itemize}

%--- Paper 18 ---
\subsubsection{Schurtenberger \& Fehr (2018): Normative Foundations}

\paragraph{Citation.}
Schurtenberger, I. \& Fehr, E. (2018).
``Normative Foundations of Human Cooperation.''
\emph{Nature Human Behaviour}.

\paragraph{Framework Integration.}
\begin{itemize}
  \item \textbf{Norms as foundation:} $\Psi_{norm}$ enables cooperation
  \item \textbf{Beyond preferences:} Norms coordinate expectations
  \item \textbf{Enforcement:} Norm violation triggers punishment
\end{itemize}

%--- Paper 19 ---
\subsubsection{Aghion, Fehr et al. (2018): Bounded Rationality and Implementation}

\paragraph{Citation.}
Aghion, P., Fehr, E., Holden, R., \& Wilkening, T. (2018).
``The Role of Bounded Rationality and Imperfect Information in Subgame Perfect Implementation.''
\emph{Journal of the European Economic Association}.

\paragraph{Framework Integration.}
\begin{itemize}
  \item \textbf{Bounded rationality:} $C^{perceived} \neq C^{actual}$
  \item \textbf{Implementation:} Mechanism design under cognitive constraints
  \item \textbf{Robustness:} Mechanisms must work for boundedly rational agents
\end{itemize}

%--- Paper 20 ---
\subsubsection{Cohn, Fehr \& Mar\'echal (2017): Banking Culture and Risk}

\paragraph{Citation.}
Cohn, A., Fehr, E., \& Mar\'echal, M.A. (2017).
``Do Professional Norms in the Banking Industry Favor Risk-Taking?''
\emph{Review of Financial Studies}.

\paragraph{Core Finding.}
Making banker identity salient increases dishonesty. Professional norms
in banking may encourage problematic behavior.

\paragraph{Framework Integration.}
\begin{itemize}
  \item \textbf{Identity priming:} Activating $U^{IDN}_{banker}$ changes behavior
  \item \textbf{Corporate culture:} $\Psi_{inst}^{bank}$ shapes individual decisions
  \item \textbf{Financial crisis link:} Culture contributed to 2008 crisis
\end{itemize}

\subsection{Synthesis: Fehr's Contribution to the Framework}

\paragraph{Core Themes Across 20 Papers.}

\begin{enumerate}
  \item \textbf{Heterogeneity:} People differ systematically in social preferences
  \item \textbf{Stability:} Preference types are stable and predictive
  \item \textbf{Context-dependence:} Behavior varies with $\Psi$ (norms, institutions, culture)
  \item \textbf{Complementarity:} Mechanisms interact (trust $\times$ contracts, repetition $\times$ competition)
  \item \textbf{Cross-domain prediction:} Lab measures predict field behavior
  \item \textbf{Policy relevance:} Understanding preferences enables better policy
\end{enumerate}

\paragraph{Framework Elements Established.}

\begin{center}
\renewcommand{\arraystretch}{1.2}
\small
\begin{tabular}{lll}
\hline
\textbf{Element} & \textbf{Papers} & \textbf{Contribution} \\
\hline
$C^{SS}_{ij} \neq 0$ & All & Core finding: social preferences exist \\
Preference types & 1, 5, 15 & Clustering reveals discrete types \\
$\Psi_{norm}$ & 16, 18 & Norms as coordination device \\
$\Psi_{culture}$ & 20 & Professional culture affects decisions \\
$C^{mechanism} > 0$ & 4, 9 & Mechanisms are complements \\
Lab $\to$ Field & 3, 7, 13 & Experimental validity confirmed \\
Policy design & 2, 8, 11, 19 & Behavioral constraints on mechanisms \\
\hline
\end{tabular}
\end{center}

\subsection{Bonus: Female Genital Cutting---Norm Change in Practice}

Three additional papers by Fehr and collaborators address female genital
cutting (FGC), demonstrating how behavioral economics can inform policy
on deeply entrenched cultural practices.

%--- FGC Paper 1 ---
\subsubsection{Efferson et al. (2015): FGC Is Not a Coordination Norm}

\paragraph{Citation.}
Efferson, C., Vogt, S., Elhadi, A., Ahmed, H.E.F., \& Fehr, E. (2015).
``Female genital cutting is not a social coordination norm.''
\emph{Science}, 349(6255): 1446-1447.

\paragraph{Core Finding.}
Contrary to influential theories, FGC is \emph{not} maintained primarily
by marriage market coordination. The practice persists even when
marriageability is not at stake.

\paragraph{Framework Integration.}
\begin{itemize}
  \item \textbf{Norm misdiagnosis:} Standard $\Psi_{norm}^{coordination}$ model fails
  \item \textbf{Alternative mechanisms:} Identity ($U^{IDN}$), in-group signaling
  \item \textbf{Policy implication:} Coordination-based interventions may fail
\end{itemize}

%--- FGC Paper 2 ---
\subsubsection{Vogt et al. (2016): Changing Cultural Attitudes}

\paragraph{Citation.}
Vogt, S., Zaid, N.A.M., Ahmed, H.E.F., Fehr, E., \& Efferson, C. (2016).
``Changing cultural attitudes towards female genital cutting.''
\emph{Nature}, 538(7626): 506-509.

\paragraph{Core Finding.}
Entertainment-education movies highlighting \emph{within-community}
disagreement can shift attitudes toward FGC. External pressure is
less effective than revealing internal dissent.

\paragraph{Framework Integration.}
\begin{itemize}
  \item \textbf{Information intervention:} $\Delta\Psi_{information} \to \Delta$ attitudes
  \item \textbf{Endogenous norms:} Revealing heterogeneity destabilizes $\Psi_{norm}$
  \item \textbf{Key insight:} Change from within, not external imposition
\end{itemize}

%--- FGC Paper 3 ---
\subsubsection{Vogt, Efferson \& Fehr (2017): FGC Risk in Europe}

\paragraph{Citation.}
Vogt, S., Efferson, C., \& Fehr, E. (2017).
``The risk of female genital cutting in Europe: Comparing immigrant attitudes
toward uncut girls with attitudes in a practicing country.''
\emph{SSM-Population Health}, 3: 283-293.

\paragraph{Core Finding.}
Immigrant attitudes toward FGC differ significantly from origin-country
attitudes, suggesting acculturation effects on $\Psi_{norm}$.

\paragraph{Framework Integration.}
\begin{itemize}
  \item \textbf{$\Psi$ migration:} Norms change with context transition
  \item \textbf{Acculturation:} $\Psi_{norm,t+1} = f(\Psi_{origin}, \Psi_{host}, t)$
  \item \textbf{Policy:} Risk assessment must account for norm evolution
\end{itemize}

\paragraph{Synthesis: FGC as Framework Application.}
These three papers demonstrate how the framework applies to a challenging
real-world problem:
\begin{itemize}
  \item Standard coordination model ($C^*$ based on marriage) is wrong
  \item Actual mechanisms involve identity, signaling, heterogeneity
  \item Interventions must target correct $\Psi$ components
  \item Norms evolve with context---change is possible
\end{itemize}

This is behavioral economics at its most consequential: understanding
mechanisms to design effective interventions for practices that harm
millions of girls and women.

\subsection{The Zurich Tradition}

Ernst Fehr's research program exemplifies the methodological approach
that underlies this framework:
\begin{itemize}
  \item \textbf{Experimental rigor:} Laboratory control with field validation
  \item \textbf{Theoretical grounding:} Formal models (Fehr-Schmidt) tested empirically
  \item \textbf{Heterogeneity:} Population distributions, not representative agents
  \item \textbf{Policy orientation:} Understanding for intervention design
\end{itemize}

This appendix demonstrates that the framework can systematically organize
a major research program spanning 15 years and 20+ papers. The recurring
themes---complementarity, context-dependence, heterogeneity, cross-domain
effects---are precisely what the $(C, \Psi)$ framework is designed to capture.

\subsection{Cross-Reference: Fehr's Methodological Position}

For a deeper analysis of Ernst Fehr's \textit{methodological} commitments---beyond
his empirical contributions catalogued above---see \textbf{Appendix XVIII (LIT-FEHR-METHOD):
Integration vs. Falsifiability}. That appendix reconstructs Fehr's implicit philosophy
of science: his insistence on falsifiable models, his suspicion of integrative frameworks,
and the tension between model integration and scientific discipline. The central insight:
complementarity ($\gamma$) is scientifically legitimate when it functions as an
\textit{exclusion principle}, not an integration principle.

%%%%%%%%%%%%%%%%%%%%%%%%%%%%%%%%%%%%%%%%%%%%%%%%%%%%%%%%%%%%%%%%%%%%%%%%%%%%%%%
